\section{群内計画(Within Design)の理論}
\label{b20_WithinDesign}

\subsection{授業内容}
\subsubsection{概要}
これまで学んだ群間計画に続き，本コマでは群内計画(Within-subjects Design)について学ぶ。群内計画は，同一の実験参加者が複数の条件を経験する実験デザインであり，事前-事後測定や反復測定などがこれに該当する。群内計画の最大の特徴は，個人差を個別に識別し，誤差から分離できることである。このことにより分析の精度が高まり，効果の検出力が向上する。本コマでは，群内計画の基本的な考え方，データの構造，効果と誤差の分離方法，個人差の取り扱い方について理解し，さらに階層モデルとしての統計的特性についても学ぶ。

\subsubsection{コマ主題細目}
\begin{description}
	\item[群内計画の基本構造] 群内計画(Within-subjects Design)は，同一の実験参加者が複数の条件を経験する実験デザインである。このデザインの特徴，利点と欠点，適用場面について理解する。群内計画の代表的な例として，事前-事後測定，反復測定，対応のある比較などがある。個人差を誤差から分離できるという最大の利点がある一方で，練習効果，疲労効果，順序効果といった問題も生じうることを理解する。また，実験のコストや実施上の制約，倫理的配慮についても考察する。群内計画を表記する際の規則(例：「3水準の一要因Within計画」)についても学ぶ。
	      \begin{flushright}
		      $\to$ \textcite{Toyoda201706}のP.45--62，\textcite{kosugi2018}のP.74--78.
	      \end{flushright}

	\item[個人差の分離と精度の向上] 群内計画の最大の特徴は，個人差を誤差から分離できることである。従来の群間計画では，個人差は誤差に含まれていたが，群内計画では個人要因として明示的に扱うことができる。これにより比較すべき誤差(残差)が小さくなり，効果の検出力が向上する。個人差の分散と残差の分散を分離する方法，それによってF値がどのように変化するかを理解する。また，効果量の算出方法と解釈についても学ぶ。視覚的な例を通じて，個人差を分離することの意義を直感的に理解する。
	      \begin{flushright}
		      $\to$ \textcite{Yamada2004}のP.178--183，\textcite{kosugi2019}のP.146--148.
	      \end{flushright}

	\item[群内計画のモデルと平方和の分解] 群内計画のデータ構造を数学的に表現し，平方和の分解方法を理解する。総平方和($SS_T$)は，要因の効果による平方和($SS_A$)，個人差による平方和($SS_S$)，および残差による平方和($SS_E$)に分解される($SS_T = SS_A + SS_S + SS_E$)。各平方和の計算方法と意味を理解し，効果と誤差，個人差をどのように分離するかを学ぶ。また，適切な比較対象となる平方和(誤差項)の選択方法と，F値の算出方法についても理解する。群内計画の分散分析表の読み方と，結果の解釈方法についても学ぶ。
	      \begin{flushright}
		      $\to$ \textcite{Yamada2004}のP.179，図7.5.1の平方和の分解図式，\textcite{kosugi2018}のP.81--88.
	      \end{flushright}

	\item[群内計画と階層モデル] 群内計画は，統計的には階層モデル(multilevel model)あるいは混合モデル(mixed model)として捉えることができる。個人差という別の分布が入った統計モデルとして理解することで，より複雑な実験デザインにも対応できる柔軟性を持つ。固定効果(fixed effect)と変量効果(random effect)の違い，個人要因を変量効果として扱う意義，変量効果モデルの特性について学ぶ。また，個人差が正規分布に従うという仮定の妥当性と，その検証方法についても理解する。階層モデルの考え方は，より複雑な統計モデル(一般化線形混合モデルなど)への橋渡しとなることについても触れる。
	      \begin{flushright}
		      $\to$ \textcite{kosugi2018}のP.136--138，階層モデルについては\textcite{Toyoda201706}のP.121--145.
	      \end{flushright}

\end{description}
\subsubsection{キーワード}
\begin{itemize}
	\item 群内計画(Within-subjects Design)
	\item 個人差の分離と平方和の分解
	\item 反復測定と対応のある比較
	\item 階層モデルと混合モデル
	\item 固定効果と変量効果
\end{itemize}
\subsection{授業情報}
\paragraph{コマの展開方法}
講義
\paragraph{標準シラバスにおける位置づけ}
科目番号4；心理学研究法；2.データを用いた実証的な思考方法；C データの統計的記述\\
科目番号5；心理学統計法；1.心理学で用いられる統計手法；G 推測統計:代表値と散布度をめぐって(1)
\subsubsection{予習・復習課題}
\paragraph{予習}
前回までに学んだ群間計画(Between Design)の特徴と分散分析の基本的な考え方を復習しておくこと。特に，一要因の分散分析における平方和の分解，F値の意味について理解を確認しておくとよい。また，実験計画の用語(要因，水準など)についても再確認しておくこと。心理学の論文から群内計画を用いた研究例を探して読んでみることも有用である。

\paragraph{復習}
授業で学んだ内容を定着させるために，以下の点について理解を深めておくこと：

\begin{enumerate}
\item 群間計画と群内計画の違い，それぞれの利点と欠点を説明できるようにする
\item 群内計画における平方和の分解方法を理解し，個人差が誤差から分離される過程を説明できるようにする
\item 群内計画における適切なF値の算出方法を理解し，分散分析表を正しく解釈できるようにする
\item 個人差を変量効果として扱う階層モデル・混合モデルの考え方を理解し，その特徴を説明できるようにする
\item 実際の研究例をもとに，群内計画が適切な状況とそうでない状況を判断できるようにする
\end{enumerate}

群内計画における個人差を算出するための手続きを，具体的な数値例を用いて自分で確認してみるとよい。公開されている実験データや，教科書に載っている例題を使って練習するとよい。また，混合モデル(階層線形モデル，マルチレベルモデルとも呼ばれる)については，さまざまな呼称があるため，それぞれのキーワードで参考文献を探してみるとよい。